% handout.tex with one line more at the top, so every word below it stands a line lower: the copy the showcase holds after its marks were made on the first.
\documentclass[11pt]{article}
\usepackage{amsmath,amssymb}
\usepackage[margin=1.2in]{geometry}
\pagestyle{empty}

\begin{document}

\noindent\textbf{Balanced quivers: a reading handout}

\medskip
\noindent For the seminar on flow polytopes.

\noindent Revised after the first meeting.

\bigskip
A quiver is balanced when at every vertex the arrows coming in and the arrows going out carry the same total weight. The flow polytope of a balanced quiver is integral, and its lattice points count the integer flows.

The counting function of the polytope is a polynomial in the dilation factor. Its degree is the rank of the cycle lattice, which is the number of arrows less the number of vertices plus the number of connected components.

\begin{equation*}
  \dim \Pi(Q) = |A(Q)| - |V(Q)| + c(Q).
\end{equation*}

Two questions for the seminar. First, whether integrality needs balance at every vertex or only on average. Second, whether the leading coefficient of the counting polynomial is the normalised volume of the polytope.

\end{document}
